% Options for packages loaded elsewhere
\PassOptionsToPackage{unicode}{hyperref}
\PassOptionsToPackage{hyphens}{url}
%
\documentclass[
]{article}
\usepackage{amsmath,amssymb}
\usepackage{iftex}
\ifPDFTeX
  \usepackage[T1]{fontenc}
  \usepackage[utf8]{inputenc}
  \usepackage{textcomp} % provide euro and other symbols
\else % if luatex or xetex
  \usepackage{unicode-math} % this also loads fontspec
  \defaultfontfeatures{Scale=MatchLowercase}
  \defaultfontfeatures[\rmfamily]{Ligatures=TeX,Scale=1}
\fi
\usepackage{lmodern}
\ifPDFTeX\else
  % xetex/luatex font selection
  \setmainfont[]{DejaVu Serif}
  \setmonofont[]{DejaVu Sans Mono}
\fi
% Use upquote if available, for straight quotes in verbatim environments
\IfFileExists{upquote.sty}{\usepackage{upquote}}{}
\IfFileExists{microtype.sty}{% use microtype if available
  \usepackage[]{microtype}
  \UseMicrotypeSet[protrusion]{basicmath} % disable protrusion for tt fonts
}{}
\makeatletter
\@ifundefined{KOMAClassName}{% if non-KOMA class
  \IfFileExists{parskip.sty}{%
    \usepackage{parskip}
  }{% else
    \setlength{\parindent}{0pt}
    \setlength{\parskip}{6pt plus 2pt minus 1pt}}
}{% if KOMA class
  \KOMAoptions{parskip=half}}
\makeatother
\usepackage{xcolor}
\usepackage[margin=0.75in]{geometry}
\usepackage{longtable,booktabs,array}
\usepackage{calc} % for calculating minipage widths
% Correct order of tables after \paragraph or \subparagraph
\usepackage{etoolbox}
\makeatletter
\patchcmd\longtable{\par}{\if@noskipsec\mbox{}\fi\par}{}{}
\makeatother
% Allow footnotes in longtable head/foot
\IfFileExists{footnotehyper.sty}{\usepackage{footnotehyper}}{\usepackage{footnote}}
\makesavenoteenv{longtable}
\setlength{\emergencystretch}{3em} % prevent overfull lines
\providecommand{\tightlist}{%
  \setlength{\itemsep}{0pt}\setlength{\parskip}{0pt}}
\setcounter{secnumdepth}{-\maxdimen} % remove section numbering
\ifLuaTeX
  \usepackage{selnolig}  % disable illegal ligatures
\fi
\usepackage{bookmark}
\IfFileExists{xurl.sty}{\usepackage{xurl}}{} % add URL line breaks if available
\urlstyle{same}
\hypersetup{
  hidelinks,
  pdfcreator={LaTeX via pandoc}}

\author{}
\date{}

\begin{document}

{
\setcounter{tocdepth}{3}
\tableofcontents
}
\section{EX12 --- Rapport confirmatoire P7 : résultats de
test}\label{ex12-rapport-confirmatoire-p7-ruxe9sultats-de-test}

\textbf{Date de clôture : 2026-11-04}\\
\textbf{Statut de porte : PASS (12/12 contrôles)}

\subsection{1. Objet et portée}\label{objet-et-portuxe9e}

P7 est la campagne confirmatoire finale pré-enregistrée d'EX12. Elle
évalue quatre variantes HC-RBF (R0, C2, C3, C4) contre des MLP appariés
par nombre de paramètres entraînables, sur quatre fonctions synthétiques
de régression et vingt graines confirmatoires. La comparaison primaire
porte sur NMAE normalisée par l'écart-type de la cible calculé
exclusivement sur le jeu d'entraînement. La marge de non-infériorité est
+0,02 (différence candidat − MLP).

\textbf{Portée limitée :} P7 confirme uniquement les quatre tâches
synthétiques primaires indiquées ci-dessous. Les résultats Diabetes et
classification du pilote P3 restent exploratoires ; ils ne deviennent
pas des résultats confirmatoires par extrapolation. Aucun holdout
historique EX3--EX11 n'a été utilisé.

\subsection{2. Protocole gelé et accès au
test}\label{protocole-geluxe9-et-accuxe8s-au-test}

\begin{itemize}
\tightlist
\item
  Tâches primaires : \texttt{f1\_sin}, \texttt{f2\_mul},
  \texttt{f3\_square\_plus}, \texttt{f4\_sin\_product}.
\item
  20 graines confirmatoires disjointes des 10 graines pilotes.
\item
  Huit slots par tâche×graine : R0/C2/C3/C4 et MLP apparié
  M0\_MATCH\_R0/C2/C3/C4.
\item
  640 entraînements : 4 tâches × 20 graines × 8 slots.
\item
  Paramètres entraînables appariés à ±1 \% lorsque réalisable ; écart
  maximal vérifié : 0,9238 \%. Les centres fixes, buffers, normaliseurs
  et \texttt{P\_total} sont déclarés séparément.
\item
  Train : 1 200 observations ; validation : 240 par régime ; test : 1
  000 par régime. Régimes : interpolation, frontière, extrapolation et
  shift.
\item
  Sélection du checkpoint : interpolation validation uniquement. Les
  autres régimes de validation sont diagnostiques.
\item
  Les modèles, configs, scripts et manifests ont été gelés avant
  génération des tests. Le SHA-256 du gel est
  \texttt{a43e56047af7bc8496dcb94d5c2835c9bc553d8c24891949c4f42869361206bd}.
\item
  Après gel, 80 jeux de test (640 arrays) ont été générés et contrôlés ;
  SHA-256 du manifeste :
  \texttt{b9944b13b909fe74e9b3480f926bdd1c56950080d2431516e821e241b1dcff99}.
\item
  Scoring unique : 640 modèles × 4 régimes = 2 560 évaluations. Aucune
  ré-optimisation, sélection ou ré-entraînement post-test.
\end{itemize}

\subsection{3. Intégrité et
réconciliation}\label{intuxe9grituxe9-et-ruxe9conciliation}

\begin{longtable}[]{@{}lr@{}}
\toprule\noalign{}
Contrôle & Résultat \\
\midrule\noalign{}
\endhead
\bottomrule\noalign{}
\endlastfoot
Entraînements confirmatoires & 640/640 \\
Modèles scorés & 640/640 \\
Scores modèle×régime & 2 560/2 560 \\
Jeux test & 80/80 \\
Arrays test & 640/640 \\
Comparaisons appariées par graine et régime & 1 280/1 280 \\
Contrastes primaires (4 tâches × 4 variantes) & 16/16 \\
Artefacts gelés vérifiés & 4 257/4 257 \\
Divergences de métriques / SHA après audit & 0 \\
\end{longtable}

La recomputation des métriques depuis les prédictions sauvegardées est
activée et l'audit de score passe 11/11 contrôles. Les fichiers bruts de
prédiction, métriques, scripts et registres sont livrés dans le
workspace.

\subsection{4. Résultats primaires :
interpolation}\label{ruxe9sultats-primaires-interpolation}

Différence = NMAE(HC-RBF) − NMAE(MLP apparié). Une valeur négative
favorise HC-RBF. Chaque ligne est calculée sur 20 différences appariées
par graines ; IC à 95 \% apparié. Les p-values sont corrigées par Holm
sur les 16 contrastes primaires. La non-infériorité utilise la marge
candidat−MLP ≤ +0,02 et une correction Holm unilatérale.

\begin{longtable}[]{@{}
  >{\raggedright\arraybackslash}p{(\columnwidth - 12\tabcolsep) * \real{0.1071}}
  >{\raggedleft\arraybackslash}p{(\columnwidth - 12\tabcolsep) * \real{0.1429}}
  >{\raggedleft\arraybackslash}p{(\columnwidth - 12\tabcolsep) * \real{0.1429}}
  >{\raggedleft\arraybackslash}p{(\columnwidth - 12\tabcolsep) * \real{0.1429}}
  >{\raggedleft\arraybackslash}p{(\columnwidth - 12\tabcolsep) * \real{0.1429}}
  >{\raggedleft\arraybackslash}p{(\columnwidth - 12\tabcolsep) * \real{0.1429}}
  >{\centering\arraybackslash}p{(\columnwidth - 12\tabcolsep) * \real{0.1786}}@{}}
\toprule\noalign{}
\begin{minipage}[b]{\linewidth}\raggedright
Tâche
\end{minipage} & \begin{minipage}[b]{\linewidth}\raggedleft
Variante
\end{minipage} & \begin{minipage}[b]{\linewidth}\raggedleft
Différence moyenne
\end{minipage} & \begin{minipage}[b]{\linewidth}\raggedleft
IC95 \%
\end{minipage} & \begin{minipage}[b]{\linewidth}\raggedleft
p Holm bilatéral
\end{minipage} & \begin{minipage}[b]{\linewidth}\raggedleft
p Holm NI
\end{minipage} & \begin{minipage}[b]{\linewidth}\centering
NI ?
\end{minipage} \\
\midrule\noalign{}
\endhead
\bottomrule\noalign{}
\endlastfoot
f1\_sin & R0 & +0.08343 & {[}+0.03606, +0.13080{]} & 0.01568 & 1 &
Non \\
f1\_sin & C2 & +0.10378 & {[}+0.05125, +0.15631{]} & 0.006193 & 1 &
Non \\
f1\_sin & C3 & +0.13418 & {[}+0.09919, +0.16917{]} & 2.554e-06 & 1 &
Non \\
f1\_sin & C4 & +0.04370 & {[}+0.00465, +0.08276{]} & 0.1511 & 1 & Non \\
f2\_mul & R0 & -0.01408 & {[}-0.04922, +0.02106{]} & 0.412 & 0.1415 &
Non \\
f2\_mul & C2 & -0.02375 & {[}-0.05665, +0.00915{]} & 0.2946 & 0.03555 &
Oui \\
f2\_mul & C3 & -0.03716 & {[}-0.07581, +0.00149{]} & 0.1757 & 0.02083 &
Oui \\
f2\_mul & C4 & -0.05679 & {[}-0.10532, -0.00827{]} & 0.1449 & 0.01464 &
Oui \\
f3\_square\_plus & R0 & -0.03482 & {[}-0.05974, -0.00991{]} & 0.06074 &
0.0009645 & Oui \\
f3\_square\_plus & C2 & -0.03419 & {[}-0.05762, -0.01076{]} & 0.05273 &
0.0006808 & Oui \\
f3\_square\_plus & C3 & -0.03390 & {[}-0.06488, -0.00292{]} & 0.1511 &
0.007812 & Oui \\
f3\_square\_plus & C4 & -0.05666 & {[}-0.08450, -0.02883{]} & 0.005069 &
9.65e-05 & Oui \\
f4\_sin\_product & R0 & -0.04361 & {[}-0.06433, -0.02289{]} & 0.003961 &
2.583e-05 & Oui \\
f4\_sin\_product & C2 & -0.04659 & {[}-0.06602, -0.02716{]} & 0.001067 &
6.104e-06 & Oui \\
f4\_sin\_product & C3 & -0.03941 & {[}-0.06599, -0.01282{]} & 0.05273 &
0.0009023 & Oui \\
f4\_sin\_product & C4 & -0.03553 & {[}-0.05025, -0.02081{]} & 0.001061 &
1.626e-06 & Oui \\
\end{longtable}

\textbf{Synthèse primaire : 11/16 contrastes passent le test de
non-infériorité ajusté. 7/16 rejettent l'égalité à zéro après correction
de Holm bilatérale.} Les contrastes significatifs ne sont pas tous à
l'avantage du même modèle : f1\_sin R0/C2/C3 défavorisent HC-RBF ;
f3\_square\_plus C4 et f4\_sin\_product R0/C2/C4 favorisent HC-RBF.

\textbf{Erratum de synthèse :} le compte exact recalculé depuis
\texttt{P7\_TEST\_PRIMARY\_HOLM\_16.csv} est 7/16 contrastes bilatéraux
significatifs après correction de Holm. Toute synthèse préliminaire
indiquant 4/16 est supersédée par ce compte audité.

\subsection{5. Généralisation par
régime}\label{guxe9nuxe9ralisation-par-ruxe9gime}

\begin{longtable}[]{@{}
  >{\raggedright\arraybackslash}p{(\columnwidth - 6\tabcolsep) * \real{0.2000}}
  >{\raggedleft\arraybackslash}p{(\columnwidth - 6\tabcolsep) * \real{0.2667}}
  >{\raggedleft\arraybackslash}p{(\columnwidth - 6\tabcolsep) * \real{0.2667}}
  >{\raggedleft\arraybackslash}p{(\columnwidth - 6\tabcolsep) * \real{0.2667}}@{}}
\toprule\noalign{}
\begin{minipage}[b]{\linewidth}\raggedright
Régime
\end{minipage} & \begin{minipage}[b]{\linewidth}\raggedleft
Contrastes tâche×variante
\end{minipage} & \begin{minipage}[b]{\linewidth}\raggedleft
Différence NMAE moyenne
\end{minipage} & \begin{minipage}[b]{\linewidth}\raggedleft
Contrastes favorables à HC-RBF
\end{minipage} \\
\midrule\noalign{}
\endhead
\bottomrule\noalign{}
\endlastfoot
interpolation & 16 & -0.00571 & 12/16 \\
boundary & 16 & -0.03192 & 15/16 \\
extrapolation & 16 & +0.20559 & 4/16 \\
shift & 16 & +0.00298 & 12/16 \\
\end{longtable}

Moyenne des 16 contrastes par régime :

\begin{longtable}[]{@{}lrrrr@{}}
\toprule\noalign{}
Tâche & Interpolation & Frontière & Extrapolation & Shift \\
\midrule\noalign{}
\endhead
\bottomrule\noalign{}
\endlastfoot
f1\_sin & +0.0913 & -0.0068 & +0.2626 & +0.1289 \\
f2\_mul & -0.0329 & -0.0257 & +0.4327 & -0.0386 \\
f3\_square\_plus & -0.0399 & -0.0353 & +0.1491 & -0.0429 \\
f4\_sin\_product & -0.0413 & -0.0599 & -0.0220 & -0.0355 \\
\end{longtable}

L'extrapolation est le point faible principal : différence moyenne
+0,20559 sur les 16 contrastes ; seulement 4/16 contrastes favorisent
HC-RBF, et ces quatre cas concernent \texttt{f4\_sin\_product}. Sur
\texttt{f1\_sin}, \texttt{f2\_mul} et \texttt{f3\_square\_plus}, les
quatre variantes sont moins bonnes que leur MLP apparié en
extrapolation. Le shift moyen est proche de zéro (+0,00298), mais le
signe dépend fortement de la tâche : \texttt{f1\_sin} est défavorable
aux sondes tandis que les trois autres tâches sont favorables. La
frontière favorise HC-RBF dans 15/16 contrastes.

\subsection{6. Coût et efficacité : confrontation
P5}\label{couxfbt-et-efficacituxe9-confrontation-p5}

Les statistiques de ressources ci-dessous proviennent du
\textbf{benchmark canonique P5 V3} et du \texttt{P5\_phase\_gate.json}
(PASS 15/15). P5 a profilé 560 modèles, 13 440 cellules timing, 1 075
200 latences élémentaires, 2 560 comparaisons appariées de latence, 320
paires mémoire et 320 paires de temps de pipeline. Le protocole
utilisait CPU mono-thread, float32, batches 1/16/64/256, deux chemins
(modèle seul et normalisation+modèle), et trois blocs indépendants de 80
mesures.

Les ratios sont candidat/MLP : \textgreater1 signifie HC-RBF plus lent
ou plus volumineux. Le chemin « full » inclut la normalisation. Le temps
d'entraînement P5 est une durée de pipeline P3/P4 incluant
initialisation, validation, diagnostics de couverture et écriture des
artefacts ; il est descriptif, pas une mesure pure de l'optimiseur. La
taille mémoire est la taille exacte des tenseurs/buffers/normaliseur;
RSS est un indicateur process-level non attribuable proprement à un
modèle. Aucune mesure d'énergie ou GPU n'est disponible.

\begin{longtable}[]{@{}
  >{\raggedright\arraybackslash}p{(\columnwidth - 16\tabcolsep) * \real{0.0857}}
  >{\raggedleft\arraybackslash}p{(\columnwidth - 16\tabcolsep) * \real{0.1143}}
  >{\raggedleft\arraybackslash}p{(\columnwidth - 16\tabcolsep) * \real{0.1143}}
  >{\raggedleft\arraybackslash}p{(\columnwidth - 16\tabcolsep) * \real{0.1143}}
  >{\raggedleft\arraybackslash}p{(\columnwidth - 16\tabcolsep) * \real{0.1143}}
  >{\raggedleft\arraybackslash}p{(\columnwidth - 16\tabcolsep) * \real{0.1143}}
  >{\raggedleft\arraybackslash}p{(\columnwidth - 16\tabcolsep) * \real{0.1143}}
  >{\raggedleft\arraybackslash}p{(\columnwidth - 16\tabcolsep) * \real{0.1143}}
  >{\raggedleft\arraybackslash}p{(\columnwidth - 16\tabcolsep) * \real{0.1143}}@{}}
\toprule\noalign{}
\begin{minipage}[b]{\linewidth}\raggedright
Variante
\end{minipage} & \begin{minipage}[b]{\linewidth}\raggedleft
p50 modèle seul B1
\end{minipage} & \begin{minipage}[b]{\linewidth}\raggedleft
p50 pipeline B1
\end{minipage} & \begin{minipage}[b]{\linewidth}\raggedleft
p50 pipeline B256
\end{minipage} & \begin{minipage}[b]{\linewidth}\raggedleft
Temps pipeline enregistré
\end{minipage} & \begin{minipage}[b]{\linewidth}\raggedleft
Stockage tenseurs
\end{minipage} & \begin{minipage}[b]{\linewidth}\raggedleft
NMAE régression P5
\end{minipage} & \begin{minipage}[b]{\linewidth}\raggedleft
Log-loss classification P5
\end{minipage} & \begin{minipage}[b]{\linewidth}\raggedleft
NMAE Diabetes P5
\end{minipage} \\
\midrule\noalign{}
\endhead
\bottomrule\noalign{}
\endlastfoot
R0 & 1.894× & 1.825× & 2.586× & 1.483× & 0.996× & 1.343× & 1.034× &
0.989× \\
C2 & 1.903× & 1.824× & 2.609× & 1.515× & 0.996× & 1.186× & 1.021× &
0.988× \\
C3 & 2.048× & 1.955× & 2.922× & 1.445× & 1.269× & 1.634× & 1.021× &
0.992× \\
C4 & 2.386× & 2.254× & 3.104× & 1.622× & 1.280× & 0.988× & 1.034× &
0.974× \\
\end{longtable}

Aucune des quatre variantes n'est plus rapide sur le batch 1 dans les
comparaisons V3 (fraction de paires plus rapides = 0 dans les 80 paires
de chaque variante, sur les deux chemins). Les ratios de latence p50
batch 1 pipeline sont compris entre 1,82× et 2,25× ; au batch 256, entre
2,59× et 3,10×. Les temps de pipeline d'entraînement sont 1,44× à 1,62×
ceux du MLP, avec la réserve méthodologique indiquée ci-dessus. Les
tailles de tenseurs sont proches pour R0/C2, mais supérieures d'environ
27--28 \% pour C3/C4. Ces mesures ne satisfont donc pas l'objectif
prioritaire d'efficacité, même lorsque certaines métriques de qualité P5
sont favorables.

Les résultats P7 de temps mural d'entraînement confirmatoire, calculés
directement depuis les 80 paires par variante, donnent en complément : -
R0 : ratio médian 3.376×, ratio moyen 3.414×, plus rapide dans 15/80
paires. - C2 : ratio médian 1.418×, ratio moyen 1.541×, plus rapide dans
19/80 paires. - C3 : ratio médian 1.534×, ratio moyen 1.728×, plus
rapide dans 8/80 paires. - C4 : ratio médian 1.661×, ratio moyen 2.005×,
plus rapide dans 3/80 paires.

Ces ratios confirmatoires sont descriptifs et ne remplacent pas le
protocole V3 P5.

\textbf{Réconciliation de comptage :} un message d'avancement antérieur
avait attribué à P5 640 modèles et 256 000 latences. Ce comptage était
incorrect et confondait une autre série d'artefacts avec le benchmark P5
canonique. Les chiffres faisant autorité sont ceux du gate P5 : 560
modèles et 1 075 200 latences. Les 640 modèles et 2 560 évaluations sont
les chiffres de P7 (campagne confirmatoire test), pas P5.

\subsection{7. Verdict scientifique}\label{verdict-scientifique}

P7 fournit un résultat mixte et fortement dépendant de la tâche :
certaines variantes HC-RBF sont non inférieures ou meilleures sur
\texttt{f2\_mul}, \texttt{f3\_square\_plus} et \texttt{f4\_sin\_product}
en interpolation, mais aucune variante n'est non inférieure sur
\texttt{f1\_sin} avec la marge pré-enregistrée. La généralisation
frontière est encourageante en moyenne ; l'extrapolation échoue
nettement sur trois des quatre fonctions.

L'objectif prioritaire d'efficacité n'est pas atteint sur la plateforme
et les tailles mesurées : R0/C2/C3/C4 sont toutes plus lentes que les
MLP appariés, demandent davantage de temps de pipeline d'entraînement et
présentent une mémoire de tenseurs plus élevée pour C3/C4. Les gains de
micro-kernel P1 (vectorisation) ne se traduisent pas en performance de
bout en bout ; le coût du modèle complet et des opérations RBF reste
déterminant.

\textbf{Conclusion :} les résultats ne valident pas HC-RBF comme
remplacement général ou plus efficace des MLP. Ils montrent des niches
de qualité prédictive sur certaines fonctions et des limites
d'extrapolation et de ressources qui réfutent toute revendication
générale de supériorité à ce stade. Le résultat est négatif pour le
critère prioritaire de rapidité/ressources, mixte pour la précision et
la généralisation.

\subsection{8. Limites}\label{limites}

\begin{enumerate}
\def\labelenumi{\arabic{enumi}.}
\tightlist
\item
  Les tests confirmatoires primaires concernent quatre fonctions
  synthétiques ; ils ne démontrent pas une généralisation à des domaines
  réels, à des réseaux profonds ou à des architectures neuronales
  générales.
\item
  Diabetes et classification ne disposent pas ici d'un test
  confirmatoire P7 ; leurs résultats P3/P5 sont exploratoires.
\item
  Les benchmarks de ressources sont CPU mono-thread sur la plateforme
  documentée ; aucun résultat GPU ou mesure énergétique n'est
  revendiqué.
\item
  La campagne de spécialistes P6 a identifié seulement sept centres
  répondant aux critères stricts stabilité+utilité, tous pour
  \texttt{g5\_sin\_mul\_plus\_square}, ce qui ne soutient pas une
  spécialisation générale.
\item
  Les lacunes de provenance P1\_RUN\_001/002 et les 123 lignes P6 v1
  sans hash de script sont décrites dans des rapports distincts. Les
  scripts P7 et les artefacts gelés sont hashés.
\end{enumerate}

\subsection{9. Artefacts sources de ce
rapport}\label{artefacts-sources-de-ce-rapport}

\begin{itemize}
\tightlist
\item
  \texttt{09\_analysis/P7/P7\_TEST\_SCORE\_AUDIT.json} --- audit 11/11.
\item
  \texttt{09\_analysis/P7/P7\_TEST\_PRIMARY\_HOLM\_16.csv} --- 16
  contrastes primaires complets.
\item
  \texttt{09\_analysis/P7/P7\_TEST\_PAIRED\_SUMMARY.csv} --- 64 résumés
  par tâche/variante/régime.
\item
  \texttt{09\_analysis/P7/P7\_TEST\_PAIRED\_BY\_SEED.csv} --- données
  appariées par graine.
\item
  \texttt{08\_benchmarks/P7/P7\_TEST\_SCORE\_REGISTRY.jsonl} ---
  registre de scoring.
\item
  \texttt{08\_benchmarks/P7/P7\_TEST\_SCORING\_SUMMARY.json} --- résumé
  du scoring.
\item
  \texttt{03\_design/P7/P7\_FINAL\_MODEL\_FREEZE.json} et
  \texttt{02\_data/P7\_confirmatory/P7\_TEST\_DATA\_MANIFEST.json} ---
  gels et manifests.
\item
  \texttt{09\_analysis/P5\_V3\_LATENCY\_OVERALL.csv} et
  \texttt{P5\_V3\_PARETO\_SCORECARD.csv} --- agrégats canoniques de
  ressources.
\item
  \texttt{04\_code/audit\_analyze\_p7\_test.py} --- audit et analyse
  statistique.
\item
  \texttt{04\_code/generate\_p7\_confirmatory\_final\_report.py} ---
  génération de ce rapport.
\end{itemize}

\textbf{Rédaction :} rapport généré à partir des registres et tableaux
d'audit EX12. \textbf{Statut de signature :} document produit par le
système d'analyse de cette session ; ce n'est pas une signature
numérique certifiée ni une expertise judiciaire.

\end{document}
